\subsection{Относительность}
%%%%%%%%%%%%%%%%%%%%%%%%%%%%%%%%%%%%%%%%%%%%%%%%%%%%
%%%%%%%%%%%%%%%%%%%%%%%%%%%%%%%%%%%%%%%%%%%%%%%%%%%%
\z{Когда лодка шла по течению реки, за время $t_1 = 10\,\text{мин}$ она проходила расстояние $S_1 = 1400\,\text{м}$. А когда лодка пошла против течения, то за время $t_2 = 50\,\text{мин}$ она прошла расстояние $S_2 = 2000\,\text{м}$. Определите скорость течения реки.}
%
{$50\,\text{м/мин}$}

%%%%%%%%%%%%%%%%%%%%%%%%%%%%%%%%%%%%%%%%%%%%%%%%%%%%
\z{Стальной шарик, движущийся горизонтально со скоростью $5v$, упруго ударяется о массивную стальную вертикальную стенку, движущуюся в ту же сторону со скоростью $v$. Определите скорость шарика после отскока.}
%
{$3v$}

%%%%%%%%%%%%%%%%%%%%%%%%%%%%%%%%%%%%%%%%%%%%%%%%%%%%
\z{Эскалатор метрополитена в течение $2\,\text{мин}$ поднимает идущего вверх пассажира. Если пассажир пойдет втрое быстрее, то время подъема составит $1\,\text{мин}$. Сколько времени займет подъем по эскалатору, если тот начнет работать вдвое быстрее, а пассажир будет просто стоять на ступеньках?}
%
{$2\,\text{мин}$}

%%%%%%%%%%%%%%%%%%%%%%%%%%%%%%%%%%%%%%%%%%%%%%%%%%%%
\z{Группа туристов, двигаясь цепочкой по обочине дороги со скоростью $3.6\,\text{км/ч}$, растянулась на $400\,\text{м}$. Замыкающий послал вперед велосипедиста к руководителю группы. Доехав до руководителя, велосипедист повернул обратно. За какое время он вернется к замыкающему? Скорость велосипедиста $9\,\text{м/с}$.}
%
{$1.5\,\text{мин}$}

%%%%%%%%%%%%%%%%%%%%%%%%%%%%%%%%%%%%%%%%%%%%%%%%%%%%
\z{Определите скорость $u$, считая скорость $v$ известной (см. рисунок). Нити нерастяжимы. Тела жёсткие. Проскальзывания в блоках нет.
		\begin{center}%
 			\includegraphics[width=0.9\textwidth]{2.2.5}%
		\end{center}%
}
%
{$3v;$ ~\puan$5v;$ ~\puan$5v;$ ~\puan$5.5v$}

%%%%%%%%%%%%%%%%%%%%%%%%%%%%%%%%%%%%%%%%%%%%%%%%%%%%
\z{Определите скорость $u$, считая скорость $v$ известной (см. рисунок). Шарнирные конструкции образуют одинаковые ромбы.
		\begin{center}%
 			\includegraphics[width=0.9\textwidth]{2.2.6}%
		\end{center}%
}
%
{$\frac{4}{3}v;$ ~\puan $\frac{5}{3}v$}

%%%%%%%%%%%%%%%%%%%%%%%%%%%%%%%%%%%%%%%%%%%%%%%%%%%%
\z{Четыре черепашки находятся в вершинах квадрата со стороной $l$. Они начинают двигаться одновременно с постоянной по модулю скоростью $v$. Каждая черепаха движется по направлению к своей соседке по часовой стрелке. Где встретятся черепашки и через какое время? Какой путь они пройдут до встречи? Найдите ускорение черепахи и радиус кривизны.}
%
{$\text{Центр},\ l,\ \frac{l}{v}$}

